% GENERATED FILE. Edit canonical lessons, then run npm run generate:books.
% canonical-source-hash: fnv1a64:9bc7b117f6411dd6
% canonical-lessons: JA-C272-sagi, JA-C272-damasu, JA-C272-odosu, JA-C272-kakusu, JA-C272-hontou

\chapter{A Trick over the Phone}
\label{ch:ja-a-trick-over-the-phone}
\hlchaptermodality{272}

\begin{chapteropening}
\textbf{By the end of this chapter:} \emph{I can say how someone was tricked and threatened.}

Complete the last lesson of chapter 272: I can say how someone was tricked and threatened.
\end{chapteropening}

\section[\texorpdfstring{\ja{さぎ} (sagi)}{さぎ (sagi)}]{\ja{さぎ} (sagi) — a fraud; a swindle}

\label{lesson:JA-C272-sagi}



\begin{quote}
Before the new one: say the Japanese for to hide (oneself), then the Japanese for to catch.
\end{quote}

\subsection*{You'll want to know: \ja{さぎ}}
\textbf{\ja{さぎ}} — \emph{sagi} — \textquotedblleft{}a fraud; a swindle\textquotedblright{}.

\begin{quote}
\textbf{\ja{年}\ja{じゅう} \ja{さぎの} \ja{でんわが} \ja{あります。}} — \emph{nenjū sagi no denwa ga arimasu} — \textquotedblleft{}Swindlers phone all year round.\textquotedblright{}
\end{quote}

\begin{cousinweb}
\textbf{\ja{さぎ}} (\emph{sagi}) is Sino-Japanese (kango): two Chinese characters read with their on'yomi, readings borrowed from Chinese. They are \emph{sa}, \textquotedblleft{}deceive\textquotedblright{} (Middle Chinese \emph{tsraeH}) and \emph{gi}, \textquotedblleft{}cheat\textquotedblright{} (\emph{khi}) (KANJIDIC2; Middle Chinese in Baxter's transcription).
\end{cousinweb}

\begin{grammarlens}[title={What you've built}]
One more word for telling what happened.
\end{grammarlens}

\subsection*{Your turn}
\begin{itemize}
\raggedright
  \item \emph{Say it:} \emph{sagi}
  \item \emph{Say it:} \emph{sagi}, once more
  \item \emph{Say it:} \emph{nenjū sagi no denwa ga arimasu} — Swindlers phone all year round.
  \item \emph{Recall:} say the Japanese for to hide (oneself), then the Japanese for to catch, then say \emph{sagi} again
\end{itemize}

\begin{tcolorbox}[breakable,colback=teal!4,colframe=teal!35!black,title={Before you move on}]
What does \ja{さぎ} mean? (\textquotedblleft{}A fraud; a swindle\textquotedblright{}.) Say it once more.
\end{tcolorbox}
\section[\texorpdfstring{\ja{だます} (damasu)}{だます (damasu)}]{\ja{だます} (damasu) — to trick; to deceive}

\label{lesson:JA-C272-damasu}



\begin{quote}
Before the new one: say the Japanese for to catch, then the Japanese for a fraud.
\end{quote}

\subsection*{You'll want to know: \ja{だます}}
\textbf{\ja{だます}} — \emph{damasu} — \textquotedblleft{}to trick; to deceive\textquotedblright{}.

\begin{quote}
\textbf{\ja{その} \ja{おとこは} \ja{ははを} \ja{だましました。}} — \emph{sono otoko wa haha o damashimashita} — \textquotedblleft{}That man tricked my mother.\textquotedblright{}
\end{quote}

Polite \textbf{\ja{だまします}}, \emph{damashimasu}; a godan verb. Past (polite): \textbf{\ja{おとこが} \ja{ちちを} \ja{だましました。}}, \emph{otoko ga chichi o damashimashita}.

\begin{cousinweb}
\textbf{\ja{だます}} (\emph{damasu}) is a native Japanese word (wago), inherited rather than borrowed from Chinese; its kanji was attached later as a kun'yomi (JMdict).
\end{cousinweb}

\begin{grammarlens}[title={What you've built}]
One more word for telling what happened.
\end{grammarlens}

\subsection*{Your turn}
\begin{itemize}
\raggedright
  \item \emph{Say it:} \emph{damasu}
  \item \emph{Say it:} \emph{damasu}, once more
  \item \emph{Say it:} \emph{sono otoko wa haha o damashimashita} — That man tricked my mother.
  \item \emph{Recall:} say the Japanese for to catch, then the Japanese for a fraud, then say \emph{damasu} again
\end{itemize}

\begin{tcolorbox}[breakable,colback=teal!4,colframe=teal!35!black,title={Before you move on}]
What does \ja{だます} mean? (\textquotedblleft{}To trick; to deceive\textquotedblright{}.) Say it once more.
\end{tcolorbox}
\section[\texorpdfstring{\ja{おどす} (odosu)}{おどす (odosu)}]{\ja{おどす} (odosu) — to threaten}

\label{lesson:JA-C272-odosu}



\begin{quote}
Before the new one: say the Japanese for a fraud, then the Japanese for to trick.
\end{quote}

\subsection*{You'll want to know: \ja{おどす}}
\textbf{\ja{おどす}} — \emph{odosu} — \textquotedblleft{}to threaten\textquotedblright{}.

\begin{quote}
\textbf{\ja{はんにんは} \ja{でんわで} \ja{ははを} \ja{おどしました。}} — \emph{hannin wa denwa de haha o odoshimashita} — \textquotedblleft{}The culprit threatened my mother on the phone.\textquotedblright{}
\end{quote}

Polite \textbf{\ja{おどします}}, \emph{odoshimasu}; a godan verb. Past (polite): \textbf{\ja{どろぼうが} \ja{おとこを} \ja{おどしました。}}, \emph{dorobō ga otoko o odoshimashita}.

\begin{cousinweb}
\textbf{\ja{おどす}} (\emph{odosu}) is a native Japanese word (wago), inherited rather than borrowed from Chinese; its kanji was attached later as a kun'yomi (JMdict).
\end{cousinweb}

\begin{grammarlens}[title={What you've built}]
One more word for telling what happened.
\end{grammarlens}

\subsection*{Your turn}
\begin{itemize}
\raggedright
  \item \emph{Say it:} \emph{odosu}
  \item \emph{Say it:} \emph{odosu}, once more
  \item \emph{Say it:} \emph{hannin wa denwa de haha o odoshimashita} — The culprit threatened my mother on the phone.
  \item \emph{Recall:} say the Japanese for a fraud, then the Japanese for to trick, then say \emph{odosu} again
\end{itemize}

\begin{tcolorbox}[breakable,colback=teal!4,colframe=teal!35!black,title={Before you move on}]
What does \ja{おどす} mean? (\textquotedblleft{}To threaten\textquotedblright{}.) Say it once more.
\end{tcolorbox}
\section[\texorpdfstring{\ja{かくす} (kakusu)}{かくす (kakusu)}]{\ja{かくす} (kakusu) — to hide (something)}

\label{lesson:JA-C272-kakusu}



\begin{quote}
Before the new one: say the Japanese for to trick, then the Japanese for to threaten.
\end{quote}

\subsection*{You'll want to know: \ja{かくす}}
\textbf{\ja{かくす}} — \emph{kakusu} — \textquotedblleft{}to hide (something)\textquotedblright{}.

\begin{quote}
\textbf{\ja{はんにんは} \ja{おかねを} \ja{やまに} \ja{かくしました。}} — \emph{hannin wa okane o yama ni kakushimashita} — \textquotedblleft{}The culprit hid the money in the hills.\textquotedblright{}
\end{quote}

Polite \textbf{\ja{かくします}}, \emph{kakushimasu}; a godan verb. Past (polite): \textbf{\ja{どろぼうが} \ja{おかねを} \ja{かくしました。}}, \emph{dorobō ga okane o kakushimashita}.

\begin{cousinweb}
\textbf{\ja{かくす}} (\emph{kakusu}) is a native Japanese word (wago), inherited rather than borrowed from Chinese; its kanji was attached later as a kun'yomi (JMdict). It is the transitive partner of \emph{kakureru}, \textquotedblleft{}to hide oneself\textquotedblright{}.
\end{cousinweb}

\begin{grammarlens}[title={What you've built}]
One more word for telling what happened.
\end{grammarlens}

\subsection*{Your turn}
\begin{itemize}
\raggedright
  \item \emph{Say it:} \emph{kakusu}
  \item \emph{Say it:} \emph{kakusu}, once more
  \item \emph{Say it:} \emph{hannin wa okane o yama ni kakushimashita} — The culprit hid the money in the hills.
  \item \emph{Recall:} say the Japanese for to trick, then the Japanese for to threaten, then say \emph{kakusu} again
\end{itemize}

\begin{tcolorbox}[breakable,colback=teal!4,colframe=teal!35!black,title={Before you move on}]
What does \ja{かくす} mean? (\textquotedblleft{}To hide (something)\textquotedblright{}.) Say it once more.
\end{tcolorbox}
\section[\texorpdfstring{\ja{ほんとう} (hontō)}{ほんとう (hontō)}]{\ja{ほんとう} (hontō) — true; the truth}

\label{lesson:JA-C272-hontou}



\begin{quote}
Before the new one: say the Japanese for to threaten, then the Japanese for to hide (something).
\end{quote}

\subsection*{You'll want to know: \ja{ほんとう}}
\textbf{\ja{ほんとう}} — \emph{hontō} — \textquotedblleft{}true; the truth\textquotedblright{}.

\begin{quote}
\textbf{\ja{それは} \ja{ほんとうですか。}} — \emph{sore wa hontō desu ka} — \textquotedblleft{}Is that really so?\textquotedblright{}
\end{quote}

A na-adjective: \ja{な} before a noun.

\begin{cousinweb}
\textbf{\ja{ほんとう}} (\emph{hontō}) is Sino-Japanese (kango): two Chinese characters read with their on'yomi, readings borrowed from Chinese. They are \emph{hon}, \textquotedblleft{}true\textquotedblright{} (Middle Chinese \emph{pwonX}) and \emph{tō}, \textquotedblleft{}right\textquotedblright{} (\emph{tang}) (KANJIDIC2; Middle Chinese in Baxter's transcription).
\end{cousinweb}

\begin{grammarlens}[title={What you've built}]
One more word for telling what happened.
\end{grammarlens}

\subsection*{Your turn}
\begin{itemize}
\raggedright
  \item \emph{Say it:} \emph{hontō}
  \item \emph{Say it:} \emph{hontō}, once more
  \item \emph{Say it:} \emph{sore wa hontō desu ka} — Is that really so?
  \item \emph{Recall:} say the Japanese for to threaten, then the Japanese for to hide (something), then say \emph{hontō} again
\end{itemize}

\begin{tcolorbox}[breakable,colback=teal!4,colframe=teal!35!black,title={Before you move on}]
What does \ja{ほんとう} mean? (\textquotedblleft{}True; the truth\textquotedblright{}.) Say it once more.
\end{tcolorbox}
